%============================================================
% CHƯƠNG 6: ĐẠO VỢ CHỒNG
%============================================================
\chapter{Đạo Vợ Chồng (The Way of Spouses)}

\begin{center}
  \textit{\color{truyenblue}``A successful marriage requires falling in love many times, always with the same person.''}\\
  \textit{(Một cuộc hôn nhân thành công đòi hỏi phải yêu đi yêu lại nhiều lần, và luôn luôn là cùng một người.)}\\[4pt]
  \small\color{truyengold}--- Mignon McLaughlin ---
\end{center}
\ngancach

\section{Bữa Cơm Không Muối}
\begin{truyen}{Bữa Cơm Không Muối}{Bao Dung}
\chuhoa{N}{ }gười vợ nấu canh quên cho muối. Chồng ăn xong không nói gì.\\
\textit{(The wife cooked soup and forgot to add salt. The husband finished eating without a word.)}

Con trai hỏi: ``Bố không thấy nhạt à?''\\
\textit{(The son asked: ``Didn't you find it bland, Dad?'')}

``Có. Nhưng mẹ đã nấu sau một ngày làm việc mệt mỏi. \textbf{Nhạt một chút không sao, miễn là bữa cơm còn đủ mặt mọi người}.''\\
\textit{(``Yes. But your mother cooked it after a tiring day of work. \textbf{A little bland is fine, as long as everyone is still at the table}.'')}

Người vợ nghe được, lặng người một hồi rồi ra bếp \textbf{nấu thêm một nồi canh mới – lần này không quên muối}.\\
\textit{(The wife heard it, was silent for a moment, then went to the kitchen and \textbf{made another pot of soup – this time she did not forget the salt}.)}
\end{truyen}
\begin{baihoc}
  Bao dung trong hôn nhân không phải là bỏ qua tất cả mà là biết đâu là điều đáng nhắc nhở và đâu là điều đáng bỏ qua. Lời nói có thể làm tổn thương không kém gì im lặng đúng lúc.
\end{baihoc}

\section{Tờ Giấy Ghi Lỗi}
\begin{truyen}{Tờ Giấy Ghi Lỗi}{Trân Trọng}
\chuhoa{N}{ }gười chồng hay cằn nhằn vợ, liệt kê đủ thứ lỗi nhỏ. Vợ không phản bác.\\
\textit{(The husband often nagged his wife, listing every small fault. The wife never argued back.)}

Một hôm ông bệnh nặng nằm viện. Bà túc trực suốt đêm, chăm sóc không nghỉ.\\
\textit{(One day the husband fell gravely ill and was hospitalised. She stayed up all night caring for him without rest.)}

Ông nhìn bà và cảm thấy xấu hổ. Ông lấy tờ giấy ghi lỗi cũ ra và \textbf{xé đi từng mảnh một trước mặt bà}.\\
\textit{(He looked at her and felt ashamed. He took out the old list of faults and \textbf{tore it into pieces one by one in front of her}.)}

``Anh đã ghi những điều không đáng. Bỏ qua đi nhé.''\\
\textit{(``I had written things that were not worth writing. Please forgive me.'')}

Bà mỉm cười: ``\textbf{Em chưa bao giờ đọc tờ giấy đó. Em chỉ đọc trái tim anh thôi}.''\\
\textit{(She smiled: ``\textbf{I never read that paper. I only ever read your heart}.'')}
\end{truyen}
\begin{baihoc}
  Trong hôn nhân, điều ta chọn để ghi nhớ và điều ta chọn để tha thứ phản ánh phẩm cách của ta. Người biết buông bỏ những điều nhỏ nhặt sẽ giữ được những điều lớn lao.
\end{baihoc}

\section{Chiếc Nhẫn Cũ}
\begin{truyen}{Chiếc Nhẫn Cũ}{Chung Thủy}
\chuhoa{S}{ }au ba mươi năm hôn nhân, người vợ thấy nhẫn cưới của mình bị trầy xước và nói đùa: ``Lấy nhẫn mới đi anh.''\\
\textit{(After thirty years of marriage, the wife noticed her wedding ring was scratched and joked: ``Let us get new rings.'')}

Chồng bà cầm chiếc nhẫn cũ lên, nhìn kỹ rồi lắc đầu.\\
\textit{(Her husband picked up the old ring, looked at it carefully, then shook his head.)}

``Những vết trầy này là ba mươi năm chúng ta đã sống cùng nhau. \textbf{Nhẫn mới không có ký ức đó}.''\\
\textit{(``These scratches are the thirty years we have lived together. \textbf{A new ring would not carry those memories}.'')}

Ông đem nhẫn đi đánh bóng lại, nhưng \textbf{giữ nguyên những vết xước sâu nhất} – ghi dấu những lần khó khăn đã vượt qua cùng nhau.\\
\textit{(He had the ring polished but \textbf{kept the deepest scratches intact} – marking the difficult times they had overcome together.)}
\end{truyen}
\begin{baihoc}
  Tình yêu lâu dài không phải là tình yêu hoàn hảo mà là tình yêu chọn ở lại qua những vết sẹo. Chung thủy đẹp không phải vì không có thử thách mà vì đã vượt qua tất cả.
\end{baihoc}

\section{Bữa Sáng Mỗi Ngày}
\begin{truyen}{Bữa Sáng Mỗi Ngày}{Yêu Thương Trong Thường Nhật}
\chuhoa{S}{ }uốt bốn mươi năm hôn nhân, người chồng dậy sớm hơn vợ mười lăm phút mỗi ngày để pha cà phê.\\
\textit{(Throughout forty years of marriage, the husband woke fifteen minutes earlier than his wife every day to brew coffee.)}

Không bao giờ ông nói ``anh yêu em'' vì ông không quen.\\
\textit{(He never said ``I love you'' because he was not accustomed to it.)}

Khi ông mất, người ta hỏi vợ ông: ``Bà có bao giờ nghi ngờ tình yêu của ông không?''\\
\textit{(When he passed away, people asked his wife: ``Did you ever doubt his love?'')}

Bà nhìn ra cửa sổ nơi chiếc ấm cà phê vẫn còn đó và nói: ``\textbf{Không. Bữa sáng mỗi ngày đã nói thay ông rồi}.''\\
\textit{(She looked at the window where the coffee pot still sat and said: ``\textbf{No. Every morning's coffee spoke for him}.'')}
\end{truyen}
\begin{baihoc}
  Tình yêu trong hôn nhân không cần lời nói hoa mỹ. Nó tồn tại trong những thói quen nhỏ bé hằng ngày, những cử chỉ lặng lẽ mà chỉ người được yêu mới nhận ra.
\end{baihoc}

\section{Căn Phòng Hai Chiếc Giường}
\begin{truyen}{Căn Phòng Hai Chiếc Giường}{Tôn Trọng}
\chuhoa{K}{ }hi vào bệnh viện thăm bố, người con thấy bố và mẹ nằm hai chiếc giường cạnh nhau, bàn tay nắm vào nhau qua khe hở giữa hai giường.\\
\textit{(When visiting his father in hospital, the son saw his parents lying in adjacent beds, hands clasped together through the gap between them.)}

Cả hai đều đã già và bệnh. Nhưng \textbf{không ai buông tay}.\\
\textit{(Both were old and ill. But \textbf{neither let go of the other's hand}.)}

Y tá kể từ đêm nhập viện đến giờ họ luôn nắm tay như thế.\\
\textit{(The nurse said they had been holding hands like that since the night of admission.)}

``Bố không muốn mẹ sợ. Mẹ không muốn bố đau một mình.'' Con trai cúi đầu, \textbf{lần đầu tiên hiểu thế nào là đạo vợ chồng}.\\
\textit{(``Dad does not want Mom to be scared. Mom does not want Dad to suffer alone.'' The son bowed his head, \textbf{understanding for the first time what the way of spouses truly means}.)}
\end{truyen}
\begin{baihoc}
  Hôn nhân đích thực không phải là sự hào hứng của những ngày đầu mà là sự nắm chặt tay nhau đến những ngày cuối cùng. Đó là hình thức tình yêu sâu sắc nhất mà con người có thể trao cho nhau.
\end{baihoc}

\section{Cuộc Tranh Luận Cuối}
\begin{truyen}{Cuộc Tranh Luận Cuối}{Nhượng Bộ}
\chuhoa{V}{ }ợ chồng cãi nhau lớn nhất từ trước tới nay về chuyện tiền bạc.\\
\textit{(A husband and wife had their biggest argument ever about money.)}

Giữa chừng, người chồng đột ngột im lặng.\\
\textit{(Mid-argument, the husband suddenly fell silent.)}

``Anh im vì anh không có lý?'' – vợ hỏi.\\
\textit{(``Are you quiet because you have no argument?'' – the wife asked.)}

``Không. Anh im vì anh chợt nhớ lại ngày cưới chúng ta. \textbf{Hôm đó anh hứa sẽ là người bạn đời của em, không phải đối thủ của em}. Bây giờ anh đang tranh luận như kẻ thù của em. Điều đó sai hơn bất cứ điều gì chúng ta đang cãi nhau.''\\
\textit{(``No. I am quiet because I just remembered our wedding day. \textbf{That day I promised to be your life partner, not your opponent}. Right now I am arguing like your enemy. That is more wrong than anything we are arguing about.'')}
\end{truyen}
\begin{baihoc}
  Trong hôn nhân, mục tiêu không phải là thắng cuộc tranh luận mà là giữ gìn mối quan hệ. Người khôn ngoan biết rằng có những cuộc chiến thắng nhưng lại làm mất đi điều quý giá hơn nhiều.
\end{baihoc}

\section{Áo Mưa}
\begin{truyen}{Áo Mưa}{Quan Tâm}
\chuhoa{M}{ }ỗi hôm chồng đi làm xa, vợ luôn nhét vào túi anh chiếc áo mưa nhỏ gọn.\\
\textit{(Every day the husband left for work far away, the wife always slipped a compact raincoat into his bag.)}

Anh chưa bao giờ cần dùng đến vì anh chạy xe có mui.\\
\textit{(He never needed it since he drove a covered vehicle.)}

Nhưng anh \textbf{cũng không bao giờ lấy ra}.\\
\textit{(But he also \textbf{never removed it}.)}

Bạn bè phát hiện và hỏi: ``Sao không bỏ ra cho nhẹ túi?''\\
\textit{(Friends noticed and asked: ``Why not take it out to lighten your bag?'')}

``\textbf{Vì mỗi lần cảm thấy túi nặng là tôi cảm thấy vợ đang ở đây cùng tôi}.''\\
\textit{(``\textbf{Because every time the bag feels heavy, I feel my wife is right here with me}.'')}
\end{truyen}
\begin{baihoc}
  Tình yêu trong hôn nhân đôi khi được đo không phải bằng những cử chỉ lãng mạn mà bằng những thứ bé nhỏ ta mang theo mỗi ngày mà không ai khác biết.
\end{baihoc}

\section{Bức Thư Không Gửi}
\begin{truyen}{Bức Thư Không Gửi}{Nói Lên Cảm Xúc}
\chuhoa{S}{ }au hai mươi năm hôn nhân, người vợ tình cờ tìm thấy trong ngăn kéo chồng một xấp thư dày.\\
\textit{(After twenty years of marriage, the wife accidentally found a thick stack of letters in her husband's drawer.)}

Tất cả đều gửi cho bà. Tất cả đều chưa bao giờ được dán tem và gửi đi.\\
\textit{(They were all addressed to her. Not one had ever been stamped and mailed.)}

Thư đầu tiên hẹn hò với bà. Thư cuối cùng viết trong tuần bà bệnh nặng.\\
\textit{(The first letter was asking her out on their first date. The last was written during the week she was gravely ill.)}

Bà đọc một mình, nước mắt chảy. Ông là người \textbf{ít khi nói nhưng chưa bao giờ ngừng viết cho bà}.\\
\textit{(She read them alone, tears flowing. He was a man who \textbf{rarely spoke but had never stopped writing to her}.)}
\end{truyen}
\begin{baihoc}
  Mỗi người có cách riêng để yêu. Đừng đánh giá tình yêu chỉ bằng cách nó được bày tỏ ra bên ngoài, mà hãy nhìn vào chiều sâu và sự bền bỉ của nó.
\end{baihoc}

\section{Bát Phở Đêm Khuya}
\begin{truyen}{Bát Phở Đêm Khuya}{Đồng Hành}
\chuhoa{L}{ }úc hai giờ sáng, vợ thức dậy vì quá thèm ăn trong thời gian mang thai.\\
\textit{(At two in the morning, the pregnant wife woke up with an intense craving.)}

Chồng đang ngủ say. Bà ngần ngại không dám đánh thức vì biết anh vừa làm ca đêm mệt.\\
\textit{(The husband was fast asleep. She hesitated to wake him knowing he had just finished an exhausting night shift.)}

Nhưng khi bà vào bếp, anh đã đứng ở đó \textbf{đang nấu nồi nước phở từ lúc nào}.\\
\textit{(But when she entered the kitchen, he was already there, \textbf{having started a pot of pho broth some time before}.)}

``Anh biết em sẽ dậy đói. Anh nấu trước cho chắc.''\\
\textit{(``I knew you would wake up hungry. I cooked it ahead to be sure.'')}

Bà ôm anh giữa bếp ấm, không nói gì thêm. \textbf{Không có từ nào đủ}.\\
\textit{(She held him in the warm kitchen, saying nothing more. \textbf{No words were sufficient}.)}
\end{truyen}
\begin{baihoc}
  Tình yêu vợ chồng đẹp nhất là biết nhau đến mức đọc được cả những điều chưa nói. Đồng hành là hiểu nhau trước cả khi người kia cầu cứu.
\end{baihoc}

\section{Giọng Nói Quen Thuộc}
\begin{truyen}{Giọng Nói Quen Thuộc}{Tuổi Già}
\chuhoa{N}{ }gười vợ bị mất trí nhớ tuổi già, không còn nhận ra con cháu.\\
\textit{(The elderly wife lost her memory and could no longer recognise her children and grandchildren.)}

Nhưng mỗi khi người chồng cất tiếng gọi tên bà, \textbf{bà đều quay lại và mỉm cười}.\\
\textit{(But every time the husband called her name in his voice, \textbf{she always turned around and smiled}.)}

Bác sĩ ngạc nhiên vì bà không nhận ra ai khác.\\
\textit{(The doctor was amazed since she recognised no one else.)}

``Giọng nói của tôi đã ở với bà sáu mươi năm,'' ông giải thích nhẹ nhàng. ``\textbf{Có những thứ ngay cả bệnh tật cũng không thể lấy đi}.''\\
\textit{(``My voice has been with her for sixty years,'' he explained gently. ``\textbf{There are things that even illness cannot take away}.'')}
\end{truyen}
\begin{baihoc}
  Tình yêu thực sự khắc sâu vào tiềm thức đến mức vượt qua ranh giới của ký ức và bệnh tật. Đó là bằng chứng cuối cùng rằng hôn nhân đích thực tạo ra mối liên kết vô hình không thể phá vỡ.
\end{baihoc}
